% !TEX TS-program = xelatex
\documentclass[11pt,a4paper]{article}

\usepackage{fontspec}
\usepackage{xeCJK}
\usepackage[margin=2cm]{geometry}
\usepackage{xcolor}
\usepackage{tcolorbox}
\usepackage{graphicx}
\usepackage{enumitem}
\usepackage{hyperref}

% Fonts
\setmainfont{Charis SIL}
\setCJKmainfont{Hiragino Sans}

% Colors
\definecolor{maincolor}{RGB}{0,102,153}
\definecolor{highlightcolor}{RGB}{204,229,255}

% Title formatting
\usepackage{titlesec}
\titleformat{\section}{\Large\bfseries\color{maincolor}}{\thesection}{1em}{}

\begin{document}

\begin{center}
{\Huge\bfseries\color{maincolor} Why Can't We Say \textit{Three Cattle}?}\\[0.5em]
{\Large Understanding the Hidden Patterns in English Grammar}\\[1em]
{\large Brett Reynolds • Humber Polytechnic \& University of Toronto}
\end{center}

\vspace{0.5em}

\section*{The Puzzle}

Have you ever wondered why English grammar feels so unpredictable with words like these?

\begin{tcolorbox}[colback=highlightcolor,colframe=maincolor,boxrule=0.5pt]
\begin{itemize}[itemsep=0pt]
\item ✓ \textit{three \textbf{books}} — perfectly natural
\item ✓ \textit{three \textbf{people}} — sounds fine
\item ? \textit{three \textbf{folks}} — feels a bit odd
\item ✗ \textit{three \textbf{cattle}} — clearly wrong
\item ✗ \textit{three \textbf{furniture}} — totally unnatural
\end{itemize}
\end{tcolorbox}

Yet we can say \textit{many cattle} and \textit{many furniture} without any problem! Why do some counting words work while others don't?

\section*{The Discovery}

My research shows that English counting properties don't break randomly—they follow a \textbf{predictable order}, like peeling layers off an onion. When words resist being counted, they lose their counting abilities in a specific sequence:

\begin{center}
\begin{tabular}{rl}
\textbf{Lost first:} & \textit{a/an,} \textit{one,} \textit{three} (precise counting) \\
\textbf{Lost next:} & \textit{several} (moderate precision) \\
\textbf{Lost last:} & \textit{many,} \textit{hundreds of} (vague amounts) \\
\end{tabular}
\end{center}

This means \textit{cattle} can still work with \textit{many} even though it fails with \textit{three.} The pattern holds across hundreds of English words, and \textit{no word breaks the rule}.

\section*{Why It Matters}

\textbf{For language learners:} Understanding this pattern explains why \textit{three furnitures} sounds wrong—and why your teacher can't just say \textit{memorize this list.} There's actually a system underneath!

\textbf{For teachers:} Instead of saying \textit{these are just exceptions,} you can explain that different counting words demand different levels of precision.

\textbf{For language science:} This shows that grammar isn't just a random collection of rules. Categories like \textit{countable} and \textit{uncountable} are held together by how our brains connect word forms with meanings—and when that connection weakens, everything unravels in a predictable way.

\section*{The Big Idea}

Grammar is like an ecosystem. Just as a forest stays balanced through natural processes (not because someone designed it that way), English counting patterns stay stable because speakers unconsciously reinforce them every time they talk. When we hear \textit{many cattle,} our brains expect other counting words to work too—and when they don't, it feels wrong. That subtle discomfort is what keeps the pattern alive across generations.

\newpage

\begin{center}
{\Huge\bfseries\color{maincolor} なぜ\textit{Three Cattle}と言えないのか?}\\[0.5em]
{\Large 英語文法に隠されたパターンを理解する}\\[1em]
{\large ブレット・レイノルズ • ハンバー工科大学・トロント大学}
\end{center}

\vspace{0.5em}

\section*{謎}

なぜ英語の文法は、次のような言葉でこんなにも予測不可能に感じるのでしょうか?

\begin{tcolorbox}[colback=highlightcolor,colframe=maincolor,boxrule=0.5pt]
\begin{itemize}[itemsep=0pt]
\item ✓ \textit{three \textbf{books}} (3冊の本)— 完全に自然
\item ✓ \textit{three \textbf{people}} (3人)— 問題なし
\item ? \textit{three \textbf{folks}} (3人の人々)— 少し奇妙
\item ✗ \textit{three \textbf{cattle}} (3頭の牛)— 明らかに間違い
\item ✗ \textit{three \textbf{furniture}} (3つの家具)— 全く不自然
\end{itemize}
\end{tcolorbox}

しかし、\textit{many cattle}や\textit{many furniture}は問題なく言えます!なぜ一部の数詞は使えるのに、他は使えないのでしょうか?

\section*{発見}

私の研究では、英語の可算性の特性はランダムに壊れるのではなく、\textbf{予測可能な順序}に従うことが示されています。まるで玉ねぎの皮を剥くように、言葉が数えられることに抵抗する時、特定の順序で可算能力を失います:

\begin{center}
\begin{tabular}{rl}
\textbf{最初に失われる:} & \textit{a/an}\textit{one}\textit{three}(正確な数え方) \\
\textbf{次に失われる:} & \textit{several}(中程度の精度) \\
\textbf{最後まで残る:} & \textit{many}\textit{hundreds of}(曖昧な量) \\
\end{tabular}
\end{center}

つまり、\textit{cattle}は\textit{three}では使えなくても、\textit{many}とは使えるのです。このパターンは何百もの英単語に当てはまり、\textit{この規則を破る単語は存在しません}。

\section*{なぜ重要か}

\textbf{言語学習者にとって:} このパターンを理解することで、なぜ\textit{three furnitures}が間違って聞こえるのかが説明できます。そして、先生が単に\textit{このリストを暗記しなさい}と言えない理由も分かります。実は、その下にシステムがあるのです!

\textbf{教師にとって:} \textit{これは例外です}と言う代わりに、異なる数詞が異なるレベルの精度を要求することを説明できます。

\textbf{言語科学にとって:} これは、文法が単なるランダムな規則の集まりではないことを示しています。\textit{可算}や\textit{不可算}などのカテゴリーは、私たちの脳が単語の形と意味をどのように結びつけるかによって保たれており、その結びつきが弱まると、すべてが予測可能な方法で崩れていきます。

\section*{大きなアイデア}

文法は生態系のようなものです。森が自然のプロセスを通じてバランスを保つのと同じように(誰かがそのように設計したからではなく)、英語の数え方のパターンは、話者が話すたびに無意識に強化することで安定を保ちます。\textit{many cattle}を聞くと、私たちの脳は他の数詞も使えると期待します。そして使えない時、違和感を覚えます。その微妙な不快感こそが、世代を超えてパターンを生き続けさせているのです。

\end{document}
